
%============================================================================================================================================
% Comandos de matemática
%============================================================================================================================================

% Parte 0: Lógica proposicional.
%———————————————————————————————

% Símbolos lógicos:
\NC{\no}  {\lnot}                      % No
\NC{\Y}   {\wedge}                     % Y
\RC{\O}   {\vee}                       % O (no excluyente)
\NC{\Oex} {\veebar}                    % O excluyente
\NC{\ex}  {\exists \esp{9}}            % Existe
\NC{\exu} {\exists! \esp{9}}           % Existe un único
\NC{\nex} {\nexists \esp{9}}           % No existe
\NC{\ptd} {\esp{8} \forall \esp{8}}    % Para todo
\NC{\plt} {\esp{6} \therefore \esp{6}} % Por lo tanto
\NC{\porq}{\because}                   % Porque
\NC{\prop}{\propto}                    % Proporcional
\NC{\QED} {\blacksquare}               % Queda demostrado
\NC{\pors}{\dots}                      % Puntos suspensivos
\NC{\porv}{\vdots}                     % Puntos verticales
\NC{\porh}{\cdots}                     % Puntos centrados
\NC{\pord}{\ddots}                     % Puntos diagonales
\NC{\grs} {^{\circ}}                   % Grados
\NC{\rel} [1][]{\esp{2}             \cal{R}_{#1}\esp{3}} % Relacionado
\NC{\nrel}[1][]{\esp{-1}\not\esp{-6}\cal{R}_{#1}\esp{3}} % No relacionado

% Conjuntos numéricos:
\NC{\nin}  {\notin}           % No pertenece
\NC{\inc}  {\subseteq}        % Incluído
\NC{\ninc} {\not\subseteq}    % No incluído
\NC{\vacio}{\varnothing}      % Conjunto vacío
\NC{\compl}[1]{{#1}^{\tx{c}}} % Complemento
\NC{\union}{\cup}             % Unión
\NC{\intsc}{\cap}             % Intersección
\NC{\difs} {\vartriangle}     % Diferencia simétrica
\NC{\PC}   {\Por}             % Producto cartesiano
\RC{\inf}  {\infty}           % Infinito
\NC*{\card}[1]{                                              % Cardinal
	\begin{tikzpicture}
		\pgfmathsetlengthmacro\myWidth{.8*width("=")}            % Ancho = 80% del signo igual
		\pgfmathsetlengthmacro\myHeight{height("H")}             % Altura = altura de la letra H mayúscula
		\pgfmathsetlengthmacro\mySepY{.3333*\myWidth}            % Espacio entre líneas horizontales
		\pgfmathsetlengthmacro\mySideBearing{.1*\myWidth}        % Soporte lateral
		\def\myAngle{70}                                         % Ángulo de las líneas verticales inclinadas
		\pgfmathsetlengthmacro\mySepX{\mySepY/sin(\myAngle)}     % Calcula la separación de las líneas verticales en dirección horizontal
		\pgfmathsetlengthmacro\mySlantX{\myHeight/tan(\myAngle)} % Calcula el ancho de una línea inclinada
		\draw[line cap=round]
			(0, {(\myHeight - \mySepY)/2}) -- ++(\myWidth, 0)
			(0, {(\myHeight + \mySepY)/2}) -- ++(\myWidth, 0)
			({(\myWidth - \mySepX - \mySlantX)/2}, 0)
			-- ({(\myWidth - \mySepX + \mySlantX)/2}, \myHeight)
			({(\myWidth + \mySepX - \mySlantX)/2}, 0)
			-- ({(\myWidth + \mySepX + \mySlantX)/2}, \myHeight);
		\useasboundingbox
			(-\mySideBearing, 0)
			(\myWidth + \mySideBearing, \myHeight)
		;
	\end{tikzpicture}
	{#1}
}

% Flechas:
\NC{\To}    {\longrightarrow}     % Hacia
\NC{\ent}   {\Rightarrow}         % Entonces
\NC{\Ent}   {\Longrightarrow}     % Entonces largo
\NC{\sii}   {\Leftrightarrow}     % Si y solo si
\NC{\Sii}   {\Longleftrightarrow} % Si y solo si largo
\NC{\Tiende}{\xrightarrow}        % Tiende
\NC{\Sneg}  {\circlearrowright}   % Sentido negativo (horario)
\NC{\Spos}  {\circlearrowleft}    % Sentido positivo (anti-horario)

%————————————————————————————————————————————————————————————————————————————————————————————————————————————————————————————————————————————